\documentclass[11pt]{article}
\usepackage[T1]{fontenc}
\usepackage{lmodern}
\usepackage{amsmath,amssymb,amsthm,booktabs}
\usepackage{graphicx}
\usepackage{float}
\graphicspath{{paper/figures/}{../../paper/figures/}{../figures/}}
\usepackage[a4paper,margin=27mm]{geometry}
\usepackage[hidelinks]{hyperref}
\setlength{\emergencystretch}{3em}
% Single-file editor previews allow explicit draft placeholders.
% A build with \def\RewriteRelease{1} fails if either figure is missing.
\newif\ifrewritefiguredraft
\ifdefined\RewriteRelease\rewritefiguredraftfalse
\else\rewritefiguredrafttrue\fi
\newcommand{\missingrewritefigure}[1]{%
  \ifrewritefiguredraft
    \PackageWarning{rewrite}{Draft placeholder for missing figure #1}%
    \fbox{\parbox[c][0.22\textheight][c]{0.88\textwidth}{\centering
      DRAFT: figure file \texttt{#1} is not available in this preview.}}%
  \else
    \PackageError{rewrite}{Missing figure #1}%
      {Build from the repository root or supply the figure before release.}%
  \fi}
\newcommand{\includerewritefigure}[1]{%
  \IfFileExists{paper/figures/#1}
    {\includegraphics[width=\textwidth]{paper/figures/#1}}{%
  \IfFileExists{../figures/#1}
    {\includegraphics[width=\textwidth]{../figures/#1}}{%
  \IfFileExists{#1}
    {\includegraphics[width=\textwidth]{#1}}{\missingrewritefigure{#1}}}}}
\newtheorem{theorem}{Theorem}[section]
\newtheorem{proposition}[theorem]{Proposition}
\newtheorem{lemma}[theorem]{Lemma}
\newtheorem{corollary}[theorem]{Corollary}
\theoremstyle{remark}
\newtheorem{remark}[theorem]{Remark}
\newcommand{\R}{\mathbb R}
\newcommand{\Ranexp}{\mathbb R_{\mathrm{an},\exp}}
\newcommand{\interior}{\operatorname{int}}
\newcommand{\sinc}{\operatorname{sinc}}
\newcommand{\sinhc}{\operatorname{sinhc}}
\DeclareMathOperator{\diag}{diag}

\title{Beyond Tonnelier's conjectured bound: six cycles and uniform finiteness
in continuous piecewise-affine Li\'enard systems}
\author{Haibo Lu\\
Shanghai Institute of Technology, Shanghai 201418, People's Republic of China\\
\texttt{luhaibo1985@gmail.com}\\
ORCID: 0009-0000-2717-5968}
\date{October 3, 2026}

\begin{document}
\maketitle

\begin{abstract}
We study the continuous piecewise-affine Li\'enard system
\(\dot x=F(x)-y,\ \dot y=x\), and show that breakpoint complexity
ensures uniform finiteness but does not give Tonnelier's conjectured
upper bound.  For an explicit continuous piecewise-affine function \(F\)
with rational parameters and three genuine breakpoints, we certify at least
six distinct hyperbolic limit cycles of the corresponding system.
All four affine parts are of
focus type, and the six cycles are transverse to every switching line they
meet and persist under small parameter perturbations.  Five follow the
same full switching itinerary, exhibiting multiple fixed points within
one analytic return branch.  Independently, for each fixed \(n\), we
prove a uniform finite bound on geometrically isolated periodic orbits
for the entire family with at most \(n\) breakpoints, including grazing
cycles, with no restriction on the spectral type of the affine parts.
The proof encodes each periodic orbit
by at most \(2n+2\) affine flights and uniformly bounds each oscillatory
phase, so that the resulting family of periodic-orbit markers is definable
in \(\Ranexp\).
The resulting bound is non-effective.  Finally, adapting an
outer-breakpoint construction to the certified seed yields at least
\(n+3\) hyperbolic limit cycles for every \(n\ge3\), with all slopes
in \((-2,2)\).
\end{abstract}

\section{Introduction and main results}

We consider
\begin{equation}\label{eq:system}
 \dot x=F(x)-y,\qquad \dot y=x,
\end{equation}
where \(F:\R\to\R\) is continuous and piecewise affine.  A genuine
breakpoint is a point at which the adjacent slopes differ; it need not
be a local extremum.  A limit cycle is a geometrically isolated
periodic orbit, counted without multiplicity.  After translating \(y\),
we impose \(F(0)=0\), so the origin is the unique equilibrium.  Let
\(H_c(n)\) be the supremum of the number of
limit cycles among functions with exactly \(n\) genuine breakpoints.

Two questions should be distinguished.  The sharp counting problem asks
how many limit cycles can occur at a prescribed number of breakpoints.
The uniform finiteness problem asks whether their number is bounded
independently of the coefficients and breakpoint positions once the
number of breakpoints is fixed.

For the continuous class, Tonnelier proposed an upper bound by the
number of breakpoints \cite[Conjecture 2]{tonnelier}.  For two breakpoints,
\(H_c(2)=2\) \cite{twofold}, and the construction of
Chen--Li--Zhang--Zhang gives \(H_c(n)\ge n\) for arbitrary \(n\)
\cite{clzz}.  For \(n\ge3\), the continuous upper-bound problem is stated
as open in \cite{clzz}.  The present results
separate the two roles of breakpoint complexity: the proposed estimate
fails already at three breakpoints, whereas a bound depending only on
that number exists over the entire parameter space.

Our first main result is a robust counterexample within the all-focus
class.  Our second establishes uniform finiteness over the full parameter
space, independently of that construction.  The third propagates the
certified seed to larger breakpoint numbers.

\begin{theorem}[A robust six-cycle example]\label{thm:six}
There is a continuous piecewise-affine function with rational parameters
and exactly three genuine breakpoints whose system \eqref{eq:system} has
at least six distinct hyperbolic limit cycles.  All four affine zones are
of focus type, and the six cycles are transverse to every switching line
they meet.  Five of the six cycles follow the same full switching
itinerary.  These cycles, their switching itineraries and their stability
persist on an open neighbourhood of the parameter point in the
seven-dimensional normalised parameter space.
\end{theorem}

\begin{theorem}[Uniform finiteness at fixed complexity]\label{thm:finite}
For every fixed integer \(n\ge0\), there is a finite integer \(N_n\) such
that every system \eqref{eq:system} with a continuous piecewise-affine
function \(F\) having at most \(n\) breakpoints has at most \(N_n\)
geometrically isolated periodic orbits.  The bound is uniform in all slopes, intercepts,
breakpoint positions and orbit amplitudes.  It is non-effective in the
present proof.
\end{theorem}

\begin{theorem}[Lower bounds with more breakpoints]\label{thm:lower}
For every \(n\ge3\), there is a system with exactly \(n\) genuine
breakpoints, all slopes in \((-2,2)\), and at least \(n+3\) distinct
hyperbolic cycles transverse to all switching lines they meet.  Hence
\[
                 n+3\le H_c(n)\le N_n<\infty\qquad(n\ge3).
\]
\end{theorem}

The results distinguish two sources of cycle complexity.  A switching
itinerary records the order of events, but enumerating itineraries does
not bound the number of fixed points of an individual return branch.
The five full-itinerary cycles make this distinction explicit within a
continuous system with a unique equilibrium.  At the same time, excursion
geometry bounds the number of affine flights needed to describe a periodic
orbit, and the location of the equilibria of the affine extensions bounds
the oscillatory phase of each flight.  These bounds yield uniform
finiteness without a uniform bound on orbit amplitude or flight time.

Panazzolo's bounded-spiraling framework provides definability for flows
restricted to suitable composable cell lists
\cite[Theorem~1.2]{panazzolo}; his Section~5.5 formulates uniform finiteness
questions for semi-limit cycles in continuous piecewise-linear families.
The additional task here is to obtain a finite collection of descriptions
that covers every periodic orbit uniformly over the full Li\'enard
parameter space.  We accomplish this by bounding both the number of
flights and the oscillatory phase of each flight, including trajectories
with grazing contacts.  This yields a uniformly definable family of
periodic-orbit markers and hence a uniform bound for geometrically isolated
periodic orbits in this Li\'enard subclass.

Carmona, Fern\'andez-S\'anchez, and Novaes obtained the bound \(L^*\le8\)
for isolated crossing periodic solutions in two-zone planar piecewise-linear
systems separated by a straight line \cite{carmona}.  Bravo and
Trinidad-Forte proved finiteness for non-autonomous scalar piecewise-linear
equations on a cylinder, with bounds depending on the number of regions
and the degree of the trigonometric coefficients \cite{bravo}.  Here we
treat continuous piecewise-affine Li\'enard systems with an arbitrary
fixed number of breakpoints, including grazing periodic orbits.

Two further consequences concern the return gain and parameter dependence.
Corollary~\ref{cor:gain-oscillation} uses the alternating stability of the
certified cycles to obtain extrema of the logarithmic return gain with
alternating positive and negative values on the common branch.
Corollary~\ref{cor:parameter-cells} gives a finite partition of parameter
space with continuous ordered isolated-cycle markers and constant cycle
count on each cell; the count changes only finitely often along each
definable one-parameter path.

Section~\ref{sec:geometry} fixes the section coordinates and affine-flow
conventions.  Section~\ref{sec:six} gives the exact system, the shooting
certificate, persistence and the common analytic branch.
Section~\ref{sec:finite} encodes every periodic orbit by a bounded list of
affine flights and applies the uniform fibre theorem for \(\Ranexp\).
Section~\ref{sec:lower} adapts the outer-breakpoint construction of
Chen--Li--Zhang--Zhang \cite{clzz}, preserving hyperbolicity and
transversality in the regime \(|m_j|<2\).
Appendix~\ref{app:action} derives the action identity on the full-itinerary
branch; Appendix~\ref{app:supplement} lists the interval input.

\section{Excursion geometry and affine flight maps}\label{sec:geometry}

We first record the geometric properties used in both main proofs.  The
upper and lower axis intersections provide unique orbit markers, while the
bound on contacts with vertical lines controls the length of the later
flight coding.

For a periodic orbit write its upper and lower intersections with the
\(y\)-axis as \((0,s)\) and \((0,-r)\), respectively.  Let \(P_+(s)\)
and \(Q(r)\) denote the first return maps on the upper and lower
half-axes, respectively.  The following
elementary fact is the geometric basis for both the shooting labels and
the finite coding.

\begin{lemma}[Axis sections and vertical contacts]\label{lem:geometry}
Every nonconstant periodic orbit has exactly one intersection with each
open half of the \(y\)-axis, one positive maximum and one negative minimum
of \(x\), and at most two contacts with any vertical line, counting a
tangency once.
\end{lemma}
\begin{proof}
The orbit is a Jordan curve by uniqueness.  Since
\(\int_0^T x(t)\,dt=0\), it visits both half-planes.  Away from the
origin, the axis is crossed transversely because \(\dot x=-y\).  On an
oriented Jordan curve, transverse intersections, ordered along the
\(y\)-axis, alternate between entering and leaving the interior.
All lower-axis crossings
have the same direction, and so do all upper-axis crossings; hence each
open half-axis contains at most one intersection.  Since the orbit visits
both half-planes, crossings in both directions occur, and hence each open
half-axis is met exactly once.

At an \(x\)-critical point \(x(t_0)=b\ne0\), one has \(y(t_0)=F(b)\).
Since \(\dot x(t_0)=0\), we have \(x(t_0+\tau)-b=o(\tau)\).
Local Lipschitz continuity of \(F\), together with \(\dot y(t_0)=b\),
then gives, also when \(b\) is a breakpoint,
\[
 \dot x(t_0+\tau)=-b\tau+o(\tau),\qquad
 x(t_0+\tau)=b-\tfrac12b\tau^2+o(\tau^2).
\]
Thus \(\dot x\) crosses strictly from positive to negative at a contact
on the positive excursion and in the reverse direction on the negative
excursion.  There can be only
one of each; otherwise a second crossing would require the opposite
direction in between.  The same sign argument gives the vertical-line
bound, including at a breakpoint.
\end{proof}

For ordered breakpoints \(\xi_1<\cdots<\xi_n\) and slopes
\(m_0,\ldots,m_n\), use the continuous hinge parametrisation
\begin{equation}\label{eq:hinge}
 F_\vartheta(x)=m_0x+\sum_{j=1}^n(m_j-m_{j-1})
 \bigl((x-\xi_j)_+-(-\xi_j)_+\bigr),\qquad u_+=\max(u,0).
\end{equation}
The normalisation \(F_\vartheta(0)=0\) removes the additive constant
in \(F\), while the breakpoint positions and slopes remain free parameters.
In a zone where \(F(x)=mx+b\), the
physical flow is
\begin{equation}\label{eq:physical-flow}
 \Phi_{m,b}(t,z)=q+e^{A_mt}(z-q),\quad
 q=(0,b),\quad A_m=\begin{pmatrix}m&-1\\1&0\end{pmatrix}.
\end{equation}
For velocity \(v=F(x)-y\), \(\dot x=v\), \(\dot v=mv-x\) almost
everywhere.  In the oscillatory regime \(|m|<2\), with
\(a=m/2\), \(\omega=\sqrt{1-a^2}\),
\begin{equation}\label{eq:velocity-flow}
 \Psi_m^t=e^{at}
 \begin{pmatrix}
 \cos(\omega t)-a\sin(\omega t)/\omega&\sin(\omega t)/\omega\\
 -\sin(\omega t)/\omega&\cos(\omega t)+a\sin(\omega t)/\omega
 \end{pmatrix}.
\end{equation}
The velocity vector field can jump at a switch, while the physical vector
field is continuous; variational calculations across switches are made in
physical coordinates.

\section{A robust counterexample with three breakpoints}\label{sec:six}

We construct the six-cycle example by validating six zeros of suitable
shooting equations.  The same computation verifies their itineraries and
hyperbolicity, after which persistence follows from the implicit-function
theorem.

\subsection{The exact system and the six certified cycles}

In system \eqref{eq:system}, take \(F=F_{\vartheta_*}\) as defined by
\eqref{eq:hinge}, with
\begin{equation}\label{eq:seed}
\begin{aligned}
 (\xi_1,\xi_2,\xi_3)&=(-4.568454409748787,-1,4.697723054757359),\\
 (m_0,m_1,m_2,m_3)&=(1.7852068867006083,-0.5628649147225554,\\
 &\hspace{28mm}0.31155696868477084,-1.785205933274595).
\end{aligned}
\end{equation}
Every displayed decimal in \eqref{eq:seed} is an exact rational number.
Denote this parameter point by \(\vartheta_*\).
The adjacent slopes are distinct and lie in \((-2,2)\), so there are
three genuine breakpoints and four focus zones.

Table~\ref{tab:six} records the six certified cycles, labelled
\(\Gamma_1,\ldots,\Gamma_6\) in increasing upper-section coordinate
\(s_i^*\).  The following subsections establish their existence,
distinctness and stability.  One cycle uses the short itinerary; five
cycles traverse all four zones with the same full itinerary, specified
in \eqref{eq:charts}.  Here \(\chi\) denotes the divergence integral
over one minimal period.
\begin{table}[H]
\centering
\begin{tabular}{rlll}
\toprule
cycle&upper coordinate (approx.)&strict interval for $\chi$&stability\\\midrule
1&5.7038049296&$(-0.414496,-0.414495)$&attracting\\
2&5.9550028751&$(0.040280,0.040281)$&repelling\\
3&6.1572916585&$(-0.009958,-0.009957)$&attracting\\
4&6.5817617853&$(0.007095,0.007096)$&repelling\\
5&18.7923449781&$(-0.000234,-0.000233)$&attracting\\
6&168.8220183797&$(0.00001595,0.00001596)$&repelling\\
\bottomrule
\end{tabular}
\caption{Approximate section coordinates and rigorous enclosures of the
characteristic exponents for the six cycles.  Cycle 1 has the short
itinerary; cycles 2--6 have the full itinerary.}
\label{tab:six}
\end{table}

Figure~\ref{fig:six-return}(a) shows the short orbit
\(\Gamma_1\) and the full-itinerary orbit \(\Gamma_4\) in physical
coordinates.  All six cycles are shown in
Figure~\ref{fig:six-atlas}, Appendix~\ref{app:portraits}.

\begin{figure}[tbp]
\centering
\includerewritefigure{six-cycle-geometry-return.pdf}
\caption{Numerical trajectories and upper-section return displacement
for \eqref{eq:seed}. (a) The short orbit \(\Gamma_1\) and full orbit
\(\Gamma_4\) in physical coordinates. (b,c) The five certified full-itinerary
zeros of \(\Delta_+=P_+-s\); panel (c) uses a logarithmic horizontal
axis. Blue solid curves and filled markers denote attracting cycles;
orange dashed curves and open markers denote repelling cycles.}
\label{fig:six-return}
\end{figure}

\subsection{Shooting equations and interval validation}

For a list of prescribed nodes \(a_0,\ldots,a_k\) and slopes
\(\sigma_1,\ldots,\sigma_k\), start with \(v_0=-s\) and set
\[
 (X_j,v_j)=\Psi_{\sigma_j}^{t_j}(a_{j-1},v_{j-1}).
\]
The unknown vector is \(z=(s,t_1,\ldots,t_k)\), and the residual is
\begin{equation}\label{eq:residual}
 \mathcal R(z)=(X_1-a_1,\ldots,X_k-a_k,v_k+s).
\end{equation}
The short and full charts are
\begin{equation}\label{eq:charts}
\begin{aligned}
 k=3:&\quad(\sigma_j)=(m_2,m_1,m_2),\quad (a_j)=(0,\xi_2,\xi_2,0),\\
 k=7:&\quad(\sigma_j)=(m_2,m_1,m_0,m_1,m_2,m_3,m_2),\\
 &\quad(a_j)=(0,\xi_2,\xi_1,\xi_1,\xi_2,\xi_3,\xi_3,0).
\end{aligned}
\end{equation}

These shooting flights are split only at the prescribed nodes.  A flight
may cross the \(y\)-axis in its interior; the half-plane coding in
Section~\ref{sec:finite} subdivides it there.

\paragraph{Validation of the shooting zeros.}
We use the following contraction criterion for interval validation
\cite[Section 13]{rump}.  Let \(X\subset B\)
be nonempty closed bounded boxes in the same finite-dimensional space,
let \(z_0\in B\), and suppose that \(\mathcal R\) is \(C^1\) on a
neighbourhood of \(B\).  Let \(\mathcal J\) enclose
\(D\mathcal R(z)\) for every \(z\in B\), and let \(C\) be a fixed square
matrix.  If
\begin{equation}\label{eq:validation}
 K\subset\interior X,\qquad
 \sup_{J\in\mathcal J}\|I-CJ\|_\infty<1,
\end{equation}
where \(K\) encloses
\[
 z_0-C\mathcal R(z_0)+(I-C\mathcal J)(B-z_0),
\]
then \(\mathcal R\) has a unique zero in \(X\), lying in its interior.
Indeed, for \(T(z)=z-C\mathcal R(z)\) and \(z\in X\), the segment from
\(z_0\) to \(z\) is contained in \(B\), so the integral mean-value
formula gives \(T(z)\in K\).  The second inequality makes \(T\) a
contraction on \(X\) and makes \(C\) invertible by the Neumann series;
fixed points of \(T\) are therefore exactly zeros of \(\mathcal R\).

For cycle \(i\), the entry \(i\) of \texttt{krawczyk\_centers} in
S1's \texttt{six-cycle-001.json} is an exact rational vector
\(c_i=(\bar s_i,\bar t_{i1},\ldots,\bar t_{ik_i})\), with \(k_1=3\) and
\(k_i=7\) for \(i\ge2\).  The shooting box is
\(X_i=c_i+[-10^{-25},10^{-25}]^{k_i+1}\).  It encloses the section
coordinate and flight times, with all system parameters fixed at
\(\vartheta_*\).
Arb real-ball arithmetic \cite{arb} supplies directed enclosures at
480-bit precision.  For each exact box \(X\), an outward enclosure \(B\)
contains its rational endpoints, \(z_0\) is the exact dyadic midpoint
of \(B\), and the preconditioner \(C\) is a fixed dyadic matrix.
Automatic differentiation supplies \(\mathcal J\), and strict inclusion
is checked against the exact rational endpoints of \(X\).

\paragraph{Validation of the prescribed trajectories.}
For every flight, interval enclosures on 512 closed time cells cover the
entire normalised interval \([0,1]\).
On a monotone flight, a strict velocity sign
verifies monotonicity.  On a flight returning to the same switch, initial
and final blocks have outward and inward velocity, respectively; the
remaining cells lie strictly on the required side.  Endpoint blocks use
the exact endpoint equations, forward at departure and backward at arrival.
The final flight of the short chart starts at \(\xi_2\), crosses the
lower half-axis, turns in zone \(2\), and
reaches the upper-axis section with \(v<0\).  Its initial block has
\(v>0\); all remaining cells have \(x>\xi_2\), and every cell has
\(x<\xi_3\).  The other short turning flight also stays to the right
of \(\xi_1\), so the short orbit reaches neither outer breakpoint.
Every switching velocity is separated from zero.  The minima
\(d_{\min}\) and \(u_{\min}\) below range only over the cells where the corresponding
position or velocity inequality is imposed.  The certified bounds are
\begin{table}[H]
\centering
\begin{tabular}{lll}
\toprule
quantity&notation&strict bound\\\midrule
contraction defect&$q_{\max}$&$<7.3\cdot10^{-17}$\\
inclusion margin&$\eta_{\min}$&$>9.99\cdot10^{-26}$\\
position clearance&$d_{\min}$&$>0.0489$\\
signed velocity&$u_{\min}$&$>0.000614$\\
switching speed&$v_{\min}$&$>0.8084293584$\\
\bottomrule
\end{tabular}
\caption{Bounds used by the six-cycle interval certificate.}
\label{tab:bounds}
\end{table}

Let
\[
 z_i^*=(s_i^*,t_{i1}^*,\ldots,t_{ik_i}^*)\in\interior X_i
\]
be the unique zero certified by \eqref{eq:validation}.  Its upper-section
coordinate is recorded in Table~\ref{tab:six}.  The upper-axis boxes are
disjoint, so Lemma~\ref{lem:geometry} gives six distinct geometric orbits.

\begin{proposition}[Certified periodic-orbit data]\label{prop:certified-data}
For the exact seed \eqref{eq:seed}, each of the six boxes \(X_i\)
defined above contains a unique zero of its shooting residual.  The
corresponding flight segments satisfy the prescribed first-event and
switching-velocity inequalities, and the six upper-section coordinates are
pairwise distinct.  The first box uses the short chart and the other five use the
same full chart.
\end{proposition}

The interval construction above proves the proposition: strict inclusion
and contraction give the six zeros, while the cell checks and endpoint
blocks identify each zero with a complete trajectory.

\subsection{Hyperbolicity and persistence}

Each certified closed trajectory crosses \(x=\xi_2\) exactly twice,
and the first-event inequalities exclude any further contact with this
line.  Its primitive periodic orbit crosses the same line transversely,
so it has at least two crossings; Lemma~\ref{lem:geometry} gives at
most two.  The shooting trajectory therefore traverses one minimal
period.  In particular, the sum of its prescribed flight times is the
minimal period, and the certified divergence integral is its
characteristic exponent.

At a transverse switch the physical vector field is continuous, so the
saltation matrix is the identity.  The variational fundamental matrices
therefore compose without an additional jump factor.  In physical
coordinates, Liouville's formula and the tangent multiplier one give
\begin{equation}\label{eq:multiplier}
                       \mu=\exp(\chi).
\end{equation}
The strict intervals in Table~\ref{tab:six} avoid zero, so all six
cycles are hyperbolic.

All event equations are analytic at the certified crossings.  Their time
derivatives are the nonzero crossing velocities, so the analytic implicit
function theorem continues the flight times in the seven normalised
parameters.  If \(P_i(s,\vartheta)\) denotes the continued local return,
then at the seed
\[
 P_i(s_i^*,\vartheta_*)-s_i^*=0,\qquad
 \partial_s(P_i-s)(s_i^*,\vartheta_*)=\mu_i-1\ne0.
\]
A second application of the theorem continues the six roots in disjoint
section intervals.  Intersecting the six parameter neighbourhoods and
shrinking it if necessary keeps the slopes in \((-2,2)\), the breakpoints
ordered and genuine, and the multipliers on their original sides of one.
This proves Theorem~\ref{thm:six}.

\subsection{Coexisting cycles on one analytic return branch}

The five full-itinerary cycles lie on a common analytic return branch.
The essential point is that increasing the section amplitude increases
both excursions, so the conditions for visiting all switches hold on
one interval.

\begin{lemma}[The common full-itinerary branch]\label{lem:common-branch}
For the seed \eqref{eq:seed}, let \(J\) be the set of upper-section
amplitudes for which the first-return trajectory crosses all three
switching lines twice and transversely.  Then \(J\) is an open interval
containing \([s_2^*,s_6^*]\), and \(P_+\) is analytic on \(J\).
\end{lemma}
\begin{proof}
Put \(M=\max_j|m_j|<2\).  In \((x,v)\) coordinates the polar angle
obeys, almost everywhere,
\[
 \dot\theta=-1+\frac{m(x)xv}{x^2+v^2},\qquad
 -1-\frac M2\le\dot\theta\le-1+\frac M2<0.
\]
Every axis half-return is therefore finite.  Denote the right
half-return from \((0,-r)\) to \((0,R(r))\) in physical coordinates
by \(R(r)>0\).  Since \(\dot y=x>0\) on this half-arc, it is a
graph \(x=h_r(y)>0\) between its endpoints.  Distinct half-arcs
cannot meet by uniqueness, including at a common upper endpoint.
If \(r_2>r_1\), then
\(h_{r_2}(-r_1)>h_{r_1}(-r_1)=0\).  Were
\(R(r_2)<R(r_1)\), the ordering would reverse at \(y=R(r_2)\),
forcing an intersection.  Thus \(R(r_2)>R(r_1)\), and
\(h_{r_2}>h_{r_1}\) throughout the smaller graph's domain.  In
particular its maximal excursion \(M_R(r)=\max h_r\) is strictly
increasing.

Applying the same argument to the reflected system with
\(\widetilde F(X)=-F(-X)\), obtained by \((X,Y)=(-x,-y)\), gives a
strictly increasing left half-return \(L(s)\) from \((0,s)\) to
\((0,-L(s))\), with strictly increasing excursion
\(M_L(s)\).  These functions are continuous by continuous dependence
and the transverse axis crossings.  The upper return is
\(P_+=R\circ L\), so both \(M_L(s)\) and \(M_R(L(s))\) increase
strictly.  Each half-arc has exactly one turning point: at \(v=0\),
\(\dot v=-x\), so all such crossings have the same direction.
Consequently a negative switch is crossed twice transversely precisely
when \(M_L(s)>|\xi_i|\), and a positive switch precisely when
\(M_R(L(s))>\xi_i\).  Each condition defines an open terminal
interval; their intersection is the full-itinerary interval \(J\)
containing the five certified full orbits.  On \(J\) the six switch
hits and the axis return are transverse.  At each hit the time
derivative of the event equation is its nonzero crossing velocity.
The seven affine hitting maps, their crossing velocities and \(P_+\)
are therefore analytic by the implicit-function theorem.
\end{proof}

\begin{corollary}[Robust alternating gain extrema]
\label{cor:gain-oscillation}
There are an open neighbourhood \(V\) of \(\vartheta_*\) and an interval
\(K\subset(0,\infty)\) on which the full-itinerary upper return
\(P_+(s,\vartheta)\) is analytic for \(\vartheta\in V\).  The continued
roots \(s_2(\vartheta)<\cdots<s_6(\vartheta)\) lie in \(K\).
For every \(\vartheta\in V\), put
\(G(s,\vartheta)=\log(P_+(s,\vartheta)/s)\).  There exist points
\(u_j\in(s_{j+1}(\vartheta),s_{j+2}(\vartheta))\), \(1\le j\le4\),
such that
\begin{equation}\label{eq:alternating-critical-values}
 \partial_sG(u_j,\vartheta)=0,\qquad
                  (-1)^{j+1}G(u_j,\vartheta)>0.
\end{equation}
The critical points are chosen separately for each parameter value and
need not be nondegenerate.
\end{corollary}
\begin{proof}
By Lemma~\ref{lem:common-branch}, the full-itinerary domain \(J\) is
an open interval containing \([s_2^*,s_6^*]\).
Choose a compact interval \([a,b]\subset J\) with
these roots in its interior.  Its switch hits and axis returns are
transverse.  The analytic event equations and their first-event
conditions therefore persist on a common parameter neighbourhood by
compactness.  After shrinking this neighbourhood, they define the
full-itinerary return on \(K=(a,b)\), with all five continued roots
inside \(K\).

At each root, the return multiplier gives
\begin{equation}\label{eq:gain-root-derivative}
 \partial_sG(s_i(\vartheta),\vartheta)
 =\frac{\partial_sP_+(s_i(\vartheta),\vartheta)-1}{s_i(\vartheta)}
 =\frac{e^{\chi_i(\vartheta)}-1}{s_i(\vartheta)},\qquad 2\le i\le6.
\end{equation}
Table~\ref{tab:six} and persistence give the signs \((+,-,+,-,+)\).
For \(j=1,3\), \(G\) is positive immediately to the right of
\(s_{j+1}(\vartheta)\), while both endpoint values on
\([s_{j+1}(\vartheta),s_{j+2}(\vartheta)]\) vanish.  A positive
maximum is therefore attained at an interior critical point.
For \(j=2,4\), the same argument gives an interior negative minimum.
This proves \eqref{eq:alternating-critical-values}.  Additional zeros
between the specified roots do not affect this argument.
\end{proof}

The common itinerary fixes event order but not residence times.  If
\(t_1,\ldots,t_7\) are the full-flight times, the divergence integral
along one return is
\begin{equation}\label{eq:chi-flight}
 \chi(s)=m_2(t_1+t_5+t_7)+m_1(t_2+t_4)+m_0t_3+m_3t_6.
\end{equation}
Away from a fixed point this is a return divergence integral; at a fixed
point it is the characteristic exponent.  On \(J\), the logarithmic
return gain admits the decomposition
\[
 G(s)=\log\frac{P_+(s)}s
     =2\pi q(m_2)+A_1(s)+A_2(s)+A_3(s),
 \qquad q(m)=\frac m{\sqrt{4-m^2}}.
\]
Here \(A_i(s)\) is the net invariant jump at the two crossings of the
\(i\)-th switching line.  Appendix~\ref{app:action} defines these
actions and derives the identity and their variation directions
\((+,-,-)\) for \eqref{eq:seed}.  The action identity identifies
competing amplitude-dependent contributions along the common itinerary.
The certified roots establish that these contributions balance
repeatedly; their individual variation signs alone do not force the
observed number of roots.

Sharp counting therefore requires control of the zeros within individual
itinerary branches.  The next section obtains a uniform bound for isolated
periodic orbits through a finite encoding of their trajectories.

\section{Uniform finiteness from bounded-flight coding}\label{sec:finite}

Fix \(n\).  Use the parameter space
\[
 \mathcal P_n=\{(\xi_1,\ldots,\xi_n,m_0,\ldots,m_n):
                      \xi_1<\cdots<\xi_n\}
\]
and allow zero slope jumps; this represents systems with fewer than
\(n\) genuine breakpoints by inserting auxiliary slots.

The proof has two geometric inputs: a uniform bound on the number of
affine flights in one periodic orbit, and a uniform bound on the
oscillatory phase of each flight.  We encode these flights with
first-event conditions and prove that the codes describe exactly the
periodic orbits.  Their unique lower-axis markers then form a definable
family, to which uniform finiteness applies.

In this section, a coding flight is a segment between consecutive
contacts with the breakpoint lines or the \(y\)-axis.  Unlike the
shooting flights of Section~\ref{sec:six}, each such segment lies in a
fixed open half-plane except at its endpoints.  For the certified seed,
inserting the intermediate lower-axis crossing refines the three- and
seven-flight shooting charts into four- and eight-flight codes,
respectively.

\begin{lemma}[Uniform phase bound for half-plane flights]
\label{lem:phase-bound}
Let \(|m|<2\), \(\omega=\sqrt{1-m^2/4}\), and \(T>0\).
Suppose a nonconstant trajectory of the affine field
\(\dot x=mx+b-y\), \(\dot y=x\) remains in one open half-plane
\(x>0\) or \(x<0\) for \(0<t<T\), with both endpoints in its
closure.  Then \(\omega T\le\pi\).
\end{lemma}
\begin{proof}
The equilibrium of this affine extension is \((0,b)\).  At phase
\(\pi\), the linear flow satisfies
\[
 e^{A_m\pi/\omega}=-e^{m\pi/(2\omega)}I.
\]
If \(x(0)\ne0\), this gives
\(x(\pi/\omega)=-e^{m\pi/(2\omega)}x(0)\), which lies in the
opposite half-plane.  Thus \(T<\pi/\omega\).  If \(x(0)=0\), put
\(v_0=b-y(0)\).  Nonconstancy gives \(v_0\ne0\), and
\[
 x(t)=e^{mt/2}\frac{v_0}{\omega}\sin(\omega t).
\]
The first positive-time axis contact has phase \(\pi\), so
\(T\le\pi/\omega\) in this case as well.
\end{proof}

This is a bound on phase, not on flight time: \(\pi/\omega\) need
not remain bounded as \(|m|\to2\).  On the domains needed for these
half-plane flights, the flow in a zone has the finite spectral form
\begin{equation}\label{eq:all-flows}
e^{A_mt}=\begin{cases}
e^{at}\{\cos(\omega t)I+t\sinc(\omega t)B_m\},&m^2<4,\quad0\le\omega t\le\pi,\\
e^{at}(I+tB_m),&m^2=4,\\
e^{at}\{\cosh(\rho t)I+t\sinhc(\rho t)B_m\},&m^2>4,
\end{cases}
\end{equation}
Here \(t\ge0\), \(a=m/2\), \(B_m=A_m-aI\),
\(\omega=\sqrt{1-a^2}\) when \(|m|<2\), and
\(\rho=\sqrt{a^2-1}\) in the node case.  The quotients
\(\sinc u=\sin u/u\) and \(\sinhc u=\sinh u/u\) take the value one
at zero.  Lemma~\ref{lem:phase-bound} restricts the trigonometric
arguments to \([0,\pi]\).  Restricted analytic functions and the
exponential therefore make the resulting partial flow definable in
\(\Ranexp\), including the repeated-root and node cases.  Together
with the flight-count bound from Lemma~\ref{lem:geometry}, this is the
geometric step needed for the following proposition.

\begin{proposition}[Definability of periodic-orbit markers]
\label{prop:definable-markers}
For fixed \(n\), the set
\[
 \mathcal Z_n=\bigl\{(\vartheta,r)\in\mathcal P_n\times(0,\infty):
 (0,-r)\text{ lies on a nonconstant periodic orbit of }F_\vartheta\bigr\}
\]
is definable in \(\Ranexp\), uniformly in all slopes, ordered breakpoint
positions, and amplitudes.
\end{proposition}

We prove the proposition by encoding periodic orbits with the following
cell and first-event predicates.  The construction uses the
phase-restricted affine flows in \eqref{eq:all-flows}.
Let \(I_0=(-\infty,\xi_1)\),
\(I_j=(\xi_j,\xi_{j+1})\) for \(1\le j<n\), and
\(I_n=(\xi_n,\infty)\), with the evident convention for \(n=0\).  Put
\[
 C_{j,\varepsilon}(\vartheta)=
 \{(x,y):x\in I_j,\ \varepsilon x>0\},
 \qquad \varepsilon\in\{-1,+1\}.
\]
For fixed parameters its closure is expressed inside the same definable
structure by
\[
 z\in\overline{C_{j,\varepsilon}(\vartheta)}
 \Longleftrightarrow
 \forall\delta>0\ \exists w\in C_{j,\varepsilon}(\vartheta)
 \quad\|w-z\|^2<\delta^2.
\]
Write \(\operatorname{Flow}_j(\vartheta;t,z,w)\) for the graph of the
appropriate branch of \eqref{eq:physical-flow} and \eqref{eq:all-flows}.
Endpoint matching alone is insufficient: an affine segment may leave
its assigned zone before reaching the prescribed endpoint.  We therefore
require every intermediate point to remain in the corresponding open cell.
Define \(\operatorname{Flight}_{j,\varepsilon}(\vartheta;z,z',t)\) by
\begin{equation}\label{eq:flight-predicate}
\begin{gathered}
t>0,\quad z,z'\in\overline{C_{j,\varepsilon}(\vartheta)},\quad
\operatorname{Flow}_j(\vartheta;t,z,z'),\\
\forall u\ \bigl(0<u<t\Longrightarrow
 \exists w\,[\operatorname{Flow}_j(\vartheta;u,z,w)\ \wedge
 w\in C_{j,\varepsilon}(\vartheta)]\bigr).
\end{gathered}
\end{equation}
The final-time spectral condition makes all intermediate times admissible;
the last line is the first-event guard and ensures that no switching line
or the axis is met in the interior of a flight.

By Lemma~\ref{lem:geometry}, a periodic orbit records at most two contacts
with each breakpoint line and one upper-axis crossing.  Coincident events
are merged, so the number \(M\) of positive-duration flights satisfies
\begin{equation}\label{eq:flight-count}
                              M\le 2n+2.
\end{equation}
Let \(\mathcal B_\vartheta\) be the union of the axis and the breakpoint
lines.  Fix integers \(2\le M\le2n+2\), \(1\le h<M\), and
\(0\le j_i\le n\).  A code consists of \(r,s>0\), times \(t_i>0\), and
nodes \(z_0,\ldots,z_M\), satisfying
\begin{equation}\label{eq:code}
\begin{gathered}
r>0,\quad z_0=z_M=(0,-r),\quad z_h=(0,s),\\
z_i\in\mathcal B_\vartheta\quad(1\le i<M),\qquad
x(z_i)=0\ \Longleftrightarrow\ i=h\quad(1\le i<M),\\
\operatorname{Flight}_{j_i,\varepsilon_i}
       (\vartheta;z_{i-1},z_i,t_i)\quad(1\le i\le M),\\
\varepsilon_i=+1\ (i\le h),\qquad \varepsilon_i=-1\ (i>h).
\end{gathered}
\end{equation}
Here \(h\) marks the upper-axis crossing and \(j_i\) identifies the
zone of the \(i\)-th flight.  There are finitely many choices of these
indices.  Projecting their finite union to \((\vartheta,r)\), eliminating
the times, intermediate nodes and upper coordinate, gives a definable
set.  The next lemma identifies this projection with \(\mathcal Z_n\).

\begin{lemma}\label{lem:equivalence}
The pair \((\vartheta,r)\) admits a code \eqref{eq:code} if and only if
\((0,-r)\) lies on a nonconstant periodic orbit.
\end{lemma}
\begin{proof}
A periodic orbit supplies its contact list and hence a code by
\eqref{eq:flight-count}.  Consecutive events have positive separation;
the intervening trajectory lies in one open cell, and its oscillatory
phase, when \(|m|<2\), is bounded by Lemma~\ref{lem:phase-bound}.  The
positive excursion precedes the negative one, as prescribed in
\eqref{eq:code}.

Conversely, concatenate the flights of a code.  The first-event guards
identify each affine field with the global field on the open flight.
At the endpoints, continuity of \(F\) makes the derivatives agree with
the physical field, including at cyclic closure.  The concatenation is
therefore a \(C^1\) solution and uniqueness makes it the actual orbit.
It has positive duration and avoids the equilibrium.  There is no
interior axis contact except the single upper-axis node, so it describes
one circuit.  Projection to the unique lower-axis coordinate identifies
codes for the same orbit.
\end{proof}

\paragraph{Boundary and degenerate cases.}
A grazing contact separates two consecutive positive-duration flights
in the same cell.  A switch at zero is merged with the axis event and adds
no extra node.  Zero slope jumps merely insert auxiliary slots without
changing the field.  Empty cells contribute no flights because their
fixed-parameter closures are empty, and \(m^2=4\) uses the repeated-root
branch of \eqref{eq:all-flows}.  Centres and period annuli contribute
intervals of periodic markers, whose points are not isolated.  The
condition \(r>0\) and the open-cell guards exclude the equilibrium,
also when it lies on a switch.

\begin{lemma}[Uniform fibre-component bound]\label{lem:fibre-components}
For every definable family \(Z\subset P\times\R\) in the o-minimal
structure \(\Ranexp\), there is an integer \(B_Z\) such that every fibre
\(Z_p\) has at most \(B_Z\) connected components
\cite[Section~4.4]{vdm}.
\end{lemma}
\begin{proof}
Take a cylindrical cell decomposition compatible with \(Z\), with
the last coordinate as the fibre coordinate.  The fibre of each cell
is empty, a point, or an interval.  Hence the total number of cells
bounds the number of connected components of every fibre.
\end{proof}

\begin{proof}[Proof of Proposition~\ref{prop:definable-markers} and
Theorem~\ref{thm:finite}]
The finite definable union and projection above, together with
Lemma~\ref{lem:equivalence}, prove the marker proposition.

\emph{From isolated markers to isolated periodic orbits.}
For fixed parameters, geometric isolation of a periodic orbit is
equivalent to isolation of its lower-axis marker in
\((\mathcal Z_n)_\vartheta\).  Indeed, \(\dot x=r>0\) at
\((0,-r)\), so transversality and continuous dependence give a local
return for nearby section coordinates, including when the given orbit
grazes a switching line.  Any nearby periodic coordinate
must be a fixed point of this local return, since
Lemma~\ref{lem:geometry} excludes multiple lower-axis crossings in one
circuit.  Continuous dependence then makes the corresponding closed
orbits close.  Conversely, periodic orbits sufficiently close to the
given orbit meet its transverse section at nearby coordinates.
Thus accumulation of periodic markers and accumulation of periodic
orbits are equivalent here.  A continuous local return suffices for this
argument.

The structure \(\Ranexp\) is o-minimal
\cite[Corollary 5.13]{dmm}; Lemma~\ref{lem:fibre-components} gives a
constant \(B_n\) with
\[
 \#\pi_0((\mathcal Z_n)_\vartheta)\le B_n
 \quad\text{for every }\vartheta.
\]
Consequently the number of isolated periodic orbits is at most \(B_n\),
which proves Theorem~\ref{thm:finite}.  For \(n=0\), the normalised
linear system has no isolated cycles.
\end{proof}

\subsection{Isolated markers and parameter dependence}

\begin{corollary}[Finite parameter cells and definable scans]
\label{cor:parameter-cells}
For fixed \(n\), the parameter space \(\mathcal P_n\) has a finite
partition into definable cells \(C\) such that each \(C\) has an integer
\(k_C\ge0\) and, when \(k_C>0\), continuous definable functions
\[
 r_{C,1}<\cdots<r_{C,k_C}:C\longrightarrow(0,\infty)
\]
whose values are exactly the lower-axis markers of all geometrically
isolated periodic orbits.  In particular, their number is \(k_C\)
throughout \(C\).  Along any continuous \(\Ranexp\)-definable path
\(\gamma:T\to\mathcal P_n\), where \(T\subset\R\) is an interval,
this count is constant on a finite partition of \(T\) into points and
intervals.  Neither the parameter space nor \(T\) need be bounded.
\end{corollary}
\begin{proof}
Define the fibre-isolated marker set \(\mathcal I_n\) by
\begin{equation}\label{eq:isolated-markers}
\begin{aligned}
 (\vartheta,r)\in\mathcal I_n\quad\Longleftrightarrow\quad
 &(\vartheta,r)\in\mathcal Z_n\ \text{ and }\ \exists\varepsilon>0\\
 &\forall u>0\,
 \bigl[((\vartheta,u)\in\mathcal Z_n\ \wedge\ |u-r|<\varepsilon)
                                  \Longrightarrow u=r\bigr].
\end{aligned}
\end{equation}
It is definable by Proposition~\ref{prop:definable-markers}; the isolation
equivalence just proved identifies its fibres with the isolated cycles.
Each fibre is finite by Theorem~\ref{thm:finite}.  Apply cylindrical
cell decomposition with the marker coordinate last
\cite[Section~4.2]{vdm}, compatibly with \(\mathcal I_n\) and
\(\mathcal P_n\times\R\).  Over a parameter cell \(C\), no band cell
can belong to \(\mathcal I_n\), since it would have an interval fibre.
The remaining cells are graphs of continuous definable functions on
\(C\).  They are disjoint and hence can be ordered as displayed; they
exhaust \(\mathcal I_n\) over \(C\).  If there are none, \(k_C=0\).
Finally, each \(\gamma^{-1}(C)\) is definable in one real variable and
thus is a finite union of points and intervals.  Refining these finitely
many sets proves the assertion about the path.
\end{proof}

In particular, the isolated-cycle count changes only finitely many times
along any affine parameter path contained in \(\mathcal P_n\).
The cell decomposition is non-effective.

\section{Propagating the six-cycle construction}\label{sec:lower}

We adapt the outer-breakpoint construction of Chen--Li--Zhang--Zhang
\cite[Section 3.1]{clzz} to the certified six-cycle seed.  The construction
preserves a prescribed finite collection of hyperbolic cycles and adds an
outer cycle; a common-slope perturbation, as in the rotated-family theory
\cite{ctz}, ensures hyperbolicity and transversality of the resulting
family.  Iterating the increment from the three-breakpoint seed yields
the lower bound \(H_c(n)\ge n+3\) for \(n\ge3\).

\begin{proposition}[Outer-breakpoint increment]\label{prop:increment}
Suppose that \(F(0)=0\), that \(F\) has exactly \(q\) genuine
breakpoints, and that all its slopes lie in \((-2,2)\).  Given
\(k\ge1\) hyperbolic cycles, there is another such function with exactly
\(q+1\) genuine breakpoints, nearby continuations of the chosen cycles,
and an additional hyperbolic cycle outside them.  All \(k+1\) cycles
can be made transverse to the switching lines they meet.
\end{proposition}

\subsection{Global return and gain at infinity}

Use lower-axis coordinates and write \(D_F(r)=Q_F(r)-r\).  In velocity
coordinates \((x,v)\), with \(M=\max|m_j|<2\),
\begin{equation}\label{eq:polar}
 \dot\theta=-1+\frac{m(x)xv}{x^2+v^2}\le-1+M/2<0,\qquad
 \frac d{dt}\log\sqrt{x^2+v^2}=\frac{m(x)v^2}{x^2+v^2}.
\end{equation}
The initial point \((x,y)=(0,-r)\) is \((x,v)=(0,r)\).  Its radius
cannot vanish in finite time, and its first full clockwise turn takes
time between \(2\pi/(1+M/2)\) and \(2\pi/(1-M/2)\).  Thus the
lower return is defined for every \(r>0\).

Let \(m_-,m_+\) be the tail slopes and define
\(F_\infty(X)=m_-X\) for \(X\le0\) and \(F_\infty(X)=m_+X\)
for \(X\ge0\).  Since \(F-F_\infty\) is bounded,
\[
 \sup_X|F(rX)/r-F_\infty(X)|\le C_F/r\longrightarrow0.
\]
Under the physical scaling \((x,y)=r(X,Y)\), the vector fields have a
common Lipschitz constant.  Gronwall's inequality therefore gives
convergence of the trajectories from \((0,-1)\) on bounded time
intervals.  Scaled velocities and lifted angles converge as well, and
the strict angular bound in \eqref{eq:polar} makes the first full-turn
times converge.  The homogeneous half-plane of slope \(m\) contributes
the gain \(\exp(\pi m/\sqrt{4-m^2})\), whence
\begin{equation}\label{eq:gain}
 \lim_{r\to\infty}\frac{Q_F(r)}r
 =\rho(m_-,m_+),\quad
 \rho(u,w)=\exp\!\left(\frac{\pi u}{\sqrt{4-u^2}}+
                         \frac{\pi w}{\sqrt{4-w^2}}\right).
\end{equation}
For fixed \(u\), \(\log\rho(u,w)\) increases continuously from
\(-\infty\) to \(+\infty\) as \(w\) traverses \((-2,2)\).

\subsection{Common-slope regularisation}

\begin{lemma}\label{lem:regularization}
Let \(F_*:\R\to\R\) be continuous and piecewise affine, with
\(F_*(0)=0\), finitely many fixed breakpoints, and slopes in \((-2,2)\).
For \(F_\alpha=F_*+\alpha x\), restrict \(\alpha\) to a neighbourhood
of zero in which the slopes remain in \((-2,2)\).  The lower return
\(Q(r,\alpha)\) is defined for every \(r>0\), is jointly \(C^1\), and
satisfies \eqref{eq:rotation}.  Moreover:
\begin{enumerate}
\item If \(D(r_-,0)D(r_+,0)<0\), where
\(D(r,\alpha)=Q(r,\alpha)-r\) and \(0<r_-<r_+\), then for arbitrarily
small \(\alpha\) the interval \((r_-,r_+)\) contains a fixed point and
every fixed point in \([r_-,r_+]\) is hyperbolic.
\item Any finite collection of hyperbolic cycles at \(\alpha=0\) has
nearby hyperbolic continuations transverse to every switching line they
meet for arbitrarily small suitable \(\alpha\).
\end{enumerate}
\end{lemma}
\begin{proof}
Global return follows from \eqref{eq:polar}.

\paragraph{Step 1. \(C^1\) regularity across grazing contacts.}
On a finite non-equilibrium trajectory
segment, a switch contact is transverse or has the quadratic expansion
in Lemma~\ref{lem:geometry}; a contact on the axis is transverse.
There are consequently only finitely many contacts on the segment.
Write \(z=(x,y)\) for a reference trajectory and \(z_h\) for the
trajectory with initial-state and parameter perturbation
\(h=(\delta z_0,\delta\alpha)\).  Set \(\delta z_h=z_h-z\), and define
almost everywhere
\[
 \lambda(t)=F_*'(x(t))+\alpha,\qquad
 L(t)=\begin{pmatrix}\lambda(t)&-1\\1&0\end{pmatrix}.
\]
The remainder \(R_h\) is specified by
\[
 \delta\dot z_h=L(t)\delta z_h+\delta\alpha(x(t),0)+R_h(t),\qquad
 \delta z_h(0)=\delta z_0.
\]
Continuous dependence gives \(\delta z_h=O(\|h\|)\) uniformly on the
segment.  Let \(E_h\) be the reference times at which \(x(t)\) lies
within \(C\|h\|\) of a switch, with \(C\) large enough to cover the
state variation.  Transverse and quadratic contacts give
\(|E_h|=O(\|h\|^{1/2})\).  Outside \(E_h\), both trajectories lie in
the same affine zone, so only the product
\(\delta\alpha(x_h-x)\) contributes to \(R_h\).  The Lipschitz bound
on \(F_*\) yields
\[
 |R_h(t)|\le C\|h\|\mathbf1_{E_h}(t)+C\|h\|^2,
 \qquad \int|R_h(t)|\,dt=O(\|h\|^{3/2})=o(\|h\|).
\]
Variation of constants proves Fr\'echet differentiability.  For nearby
reference trajectories, the bounded variational coefficients converge
almost everywhere off the finite contact set, hence in \(L^1\).
Their variational solutions therefore converge uniformly, proving
continuity of the derivatives.  The terminal lower-axis crossing is
transverse, so the implicit-function theorem gives a jointly \(C^1\)
return.

\paragraph{Step 2. Strict parameter monotonicity.}
Let \(X_\alpha(x,y)=(F_*(x)+\alpha x-y,x)\).  For parameter variation
\(\eta=\partial_\alpha z\) with the initial point fixed, the wedge
\(W=\det(X_\alpha,\eta)\) satisfies
\(W'=\lambda W-x^2\) and \(W(0)=0\).  At the return time, its value is
\(-Q Q_\alpha\); variation of the return time adds a multiple of
\(X_\alpha\) and does not change the wedge.  Applying the same argument
to initial-section variation gives
\begin{equation}\label{eq:rotation}
\begin{aligned}
 Q_\alpha(r,\alpha)&=\frac1{Q(r,\alpha)}
 \int_0^T\exp\!\left(\int_t^T\lambda(u)\,du\right)x(t)^2\,dt>0,\\
 Q_r(r,\alpha)&=\frac r{Q(r,\alpha)}
 \exp\!\left(\int_0^T\lambda(t)\,dt\right)>0.
\end{aligned}
\end{equation}

\paragraph{Step 3. Hyperbolic fixed points for regular parameter values.}
Since \(D_\alpha=Q_\alpha>0\), the zero set is locally a graph
\(\alpha=A(r)\).  On the compact amplitude interval in
the first assertion, the endpoint signs persist for a sufficiently small
closed parameter interval.  The zero set in this compact rectangle is
covered by finitely many \(C^1\) graphs \(\alpha=A(r)\), and
\[
 A'(r)=-\frac{D_r(r,A(r))}{D_\alpha(r,A(r))}
\]
on each graph.  The critical values of a scalar \(C^1\) function form a
null set: on compact subintervals, cover the critical set by intervals
where \(|A'|<\varepsilon\), apply the mean value theorem, and then
exhaust the domain.  Choose \(\alpha\) outside the finite union of these
null sets.  Opposite endpoint signs give a zero; every zero in the
interval has \(D_r=Q_r-1\ne0\).  Since \(Q_r>0\) and \(Q_r\ne1\),
each such fixed point is hyperbolic.

\paragraph{Step 4. Removal of grazing along hyperbolic branches.}
For the second assertion, continue the chosen hyperbolic cycles on a
common small parameter interval.  Now \(z(t,\alpha)\) follows a periodic
orbit branch with period \(T(\alpha)\), so its parameter variation does
not have fixed initial value.  Differentiating periodicity gives
\(\eta(t+T)=\eta(t)-T'X_\alpha(t)\).  The wedge is therefore periodic
and still satisfies \(W'=\lambda W-x^2\).  With
\(\chi=\int_0^T\lambda(t)\,dt\), it follows that
\[
 W(t)=\frac{\displaystyle\int_t^{t+T}
       \exp\!\left(\int_u^{t+T}\lambda(v)\,dv\right)x(u)^2\,du}
                    {e^\chi-1}\ne0.
\]
Here hyperbolicity ensures \(e^\chi\ne1\).  At the unique positive or
negative extremum \(b(\alpha)\) of \(x\), one has
\(X_\alpha=(0,b)\), hence \(b'=-W/b\ne0\).  The derivative of this
extremal value follows from uniqueness of the extremum and the uniform
\(C^1\) expansion on a fixed phase interval.  On a sufficiently small
continuation interval its sign is fixed.  Each extremum can therefore
equal any fixed switch position at most once.  Avoiding finitely many
such parameter values makes all chosen cycles transverse.
\end{proof}

\subsection{Proof of the increment and the lower bound}

\begin{proof}[Proof of Proposition~\ref{prop:increment}]
Order the lower-section coordinates of the selected cycles as
\(0<r_1<\cdots<r_k\).  Hyperbolicity allows us to choose \(r_0>r_k\)
sufficiently close to \(r_k\) with \(D_F(r_0)\ne0\).  Choose \(R>0\)
larger than every existing breakpoint and every \(x\)-coordinate attained
on the selected cycles or on the complete return trajectory from \(r_0\).
For \(a\in(-2,2)\), different from the old right-tail slope \(m_+\), put
\[
             F_{R,a}(x)=F(x)+(a-m_+)(x-R)_+.
\]
There is exactly one new genuine breakpoint.  The selected cycles, their
local return maps, and the complete return trajectory at \(r_0\) are
unchanged; in particular, \(D_{R,a}(r_0)=D_F(r_0)\).  Choose \(a\)
so that \(\rho(m_-,a)-1\) has sign opposite to \(D_F(r_0)\).
Equation~\eqref{eq:gain} gives a finite \(r_*>r_0\) such that
\[
                  D_{R,a}(r_0)D_{R,a}(r_*)<0.
\]

Let \(\sigma_j\) be the slopes of \(F_{R,a}\) and put
\(\delta=\min_j(2-|\sigma_j|)>0\).  For \(|\alpha|<\delta/2\),
the function \(F_{R,a}+\alpha x\) keeps all slopes in \((-2,2)\)
and preserves every slope jump.  Shrinking this parameter interval
preserves both endpoint signs and continues the old cycles in disjoint
section intervals below \(r_0\).  Lemma~\ref{lem:regularization} first
provides an additional hyperbolic root in \((r_0,r_*)\); a further small
shift makes these \(k+1\) cycles transverse while preserving their
hyperbolicity and section order.  By Lemma~\ref{lem:geometry}, disjoint
periodic orbits are Jordan curves containing the origin, and their
section order makes them nested.  The additional cycle lies outside
the continuations of the chosen cycles.
\end{proof}

Starting from Theorem~\ref{thm:six}, apply the proposition \(n-3\) times.
Every step retains genuine breakpoints and slopes in \((-2,2)\), and
increases the chosen number of hyperbolic transverse cycles by one,
proving Theorem~\ref{thm:lower}.  At each step, the selected family contains
three more cycles than the number of breakpoints, giving \(6+(n-3)=n+3\).

\section{Consequences and open questions}

The breakpoint count does not give the conjectured upper bound for limit
cycles, even within the continuous all-focus class.  Five certified cycles
coexist on one analytic return branch, showing that a fixed switching
itinerary can support several distinct balances of contraction and
expansion.

Nevertheless, fixed breakpoint complexity gives a uniform finite bound
over the full parameter space.  The decisive controls are the number of
affine flights and their oscillatory phases, rather than uniform bounds
on amplitudes or periods.  The six-cycle example and the finiteness theorem
therefore describe complementary aspects of the same counting problem.
The marker theorem also gives finite variation of the isolated-cycle count
along each definable parameter path.

The results leave open the sharp value of \(H_c(3)\), effective bounds for
\(N_n\), and their growth with \(n\).  The present results give
\[
  6\le H_c(3)\le N_3<\infty,\qquad H_c(n)\ge n+3\quad(n\ge3).
\]
Analytic conditions ensuring the coexistence of several cycles within one
switching itinerary would complement the certified construction.

\appendix
\section{Action balance on the full-itinerary branch}\label{app:action}

Fix the all-focus seed \eqref{eq:seed}, and let \(J\) be the common
full-itinerary interval from Lemma~\ref{lem:common-branch}.
The origin lies in zone \(2\), so the winding contribution uses \(m_2\).

Write \(p=v=F(x)-y\) for the velocity coordinate.  In a zone of slope
\(m\), define
\[
 E_{m,x}(p)=p^2-mxp+x^2,\qquad q(m)=\frac m{\sqrt{4-m^2}}.
\]
Since \(|m|<2\), this energy is positive away from the origin.  Remove
the ray \(\{(0,p):p\le0\}\) from the \((x,p)\)-plane, and let
\(\phi_m(x,p)\) be the branch, with values in
\((-5\pi/2,-\pi/2)\), of the argument of
\[
                 \sqrt{4-m^2}\,x+i(2p-mx).
\]
The preimage of the argument cut is the same removed ray for every
\(m\).  Thus the values at both crossings of any switching line
\(x=\xi_i\ne0\) use one branch, rather than independent flightwise
lifts.  Along a flight, \(\dot x=p\) and \(\dot p=mp-x\), whence
\[
 \dot\phi_m=-\frac{\sqrt{4-m^2}}2,
 \qquad \dot E_{m,x}(p)=mE_{m,x}(p).
\]
Consequently
\begin{equation}\label{eq:focus-invariant}
 I_m(x,p)=\frac12\log E_{m,x}(p)+q(m)\phi_m(x,p),
 \qquad \partial_p I_m(x,p)=\frac p{E_{m,x}(p)},
\end{equation}
and \(I_m\) is constant on each flight.  At the upper-axis launch
\(p=-s\), the trajectory enters \(x<0\) and the one-sided angle is
\(-\pi/2\); at its return \(p=-P_+(s)\), it approaches from \(x>0\)
with angle \(-5\pi/2\).  The intervening central flight crosses
\(x=0\) with \(p>0\), away from the removed ray.

The switch \(\xi_i\) separates slopes \(m_{i-1}\) and \(m_i\).
Define
\[
                 J_i(p)=I_{m_i}(\xi_i,p)-I_{m_{i-1}}(\xi_i,p).
\]
Equation~\eqref{eq:focus-invariant} gives
\begin{equation}\label{eq:action-jump-derivative}
 J_i'(p)=\frac{\xi_i(m_i-m_{i-1})p^2}
 {E_{m_{i-1},\xi_i}(p)E_{m_i,\xi_i}(p)}.
\end{equation}
Let \(p_i^+(s)>0\) and \(p_i^-(s)<0\) be the rightward and leftward
crossing velocities, respectively, and put
\[
                    A_i(s)=J_i(p_i^+(s))-J_i(p_i^-(s)).
\]
Summing the invariant jumps over the six switch events leaves the
endpoint change in zone \(2\).  The map
\((x,p)\mapsto(\sqrt{4-m_2^2}x,2p-m_2x)\) has positive determinant
and preserves the clockwise winding.  The endpoint change is therefore
\(\log(P_+(s)/s)-2\pi q(m_2)\), so
\begin{equation}\label{eq:action-balance}
 G(s):=\log\frac{P_+(s)}s
      =2\pi q(m_2)+A_1(s)+A_2(s)+A_3(s).
\end{equation}

To determine the variation of each action, let \(T_c(s)\) be the time
of a transverse hit on a fixed switching line, and write \(y_c(s)\)
and \(p_c(s)\) for its coordinates.  The section-hit derivative is
\[
 \frac{dy_c}{ds}
 =\exp\!\left(\int_0^{T_c(s)}m(x(t;s))\,dt\right)\frac{-s}{p_c},
 \qquad
 \frac{dp_c}{ds}
 =\exp\!\left(\int_0^{T_c(s)}m(x(t;s))\,dt\right)\frac{s}{p_c}.
\]
Indeed, the flow determinant is the exponential shown, and projection
onto the hit section divides by its nonzero normal velocity \(p_c\).
The second equality uses \(p_c=F(\xi_i)-y_c\).  Thus \(p_i^+\)
increases and \(p_i^-\) decreases with \(s\).  Together with
\eqref{eq:action-jump-derivative}, this proves
\[
              \operatorname{sgn}A_i'(s)
              =\operatorname{sgn}[\xi_i(m_i-m_{i-1})].
\]
For \eqref{eq:seed} these signs are \((+,-,-)\), giving competing
amplitude-dependent contributions within the same itinerary.

For the derivative of the full return, let \(a_i\) and \(b_i\) be the
slopes in the two zones adjacent to \(\xi_i\), on the sides nearer to
and farther from the origin, respectively.
Thus \((a_i,b_i)=(m_i,m_{i-1})\) when \(\xi_i<0\), and
\((a_i,b_i)=(m_{i-1},m_i)\) when \(\xi_i>0\).  Let
\(p_i^{\rm out}\) and \(p_i^{\rm in}\) denote the velocities pointing
away from and toward the origin, and define
\[
 U_i=\frac{p_i^{\rm out}}{\xi_i}>0,
 \qquad V_i=-\frac{p_i^{\rm in}}{\xi_i}>0.
\]
These normalisations apply on both sides: the outgoing velocity is
\(p_i^-\) on the negative side and \(p_i^+\) on the positive side.
Set
\begin{equation}\label{eq:energy-ratio-appendix}
 Q_i(s)=\frac{(U_i^2-a_iU_i+1)(V_i^2+b_iV_i+1)}
 {(U_i^2-b_iU_i+1)(V_i^2+a_iV_i+1)},
\end{equation}
which is positive.  With \(T(s)\) denoting the total return time, write
\(\chi(s)=\int_0^{T(s)}m(x(t;s))\,dt\), as in
\eqref{eq:chi-flight}.  Integrating \(\dot E=mE\) on each of the
seven flights and multiplying the resulting energy ratios cancels the
intermediate endpoint energies and gives
\begin{equation}\label{eq:energy-product-appendix}
 e^{\chi(s)}=\left(\frac{P_+(s)}s\right)^2\prod_{i=1}^3Q_i(s).
\end{equation}
The section-hit derivatives likewise give
\(P_+'(s)=e^{\chi(s)}s/P_+(s)\).  Hence, with
\(B(s)=\sum_i\log Q_i(s)\) and \(\Delta_+(s)=P_+(s)-s\),
\begin{equation}\label{eq:distortion-appendix}
 e^{B(s)}=1+sG'(s),\qquad
 \Delta_+'(s)=e^{G(s)+B(s)}-1.
\end{equation}
Thus \(B=0\) locates critical points of \(G\), rather than of
\(\Delta_+\) away from fixed points.  At a fixed point \(G=0\), the
signs of \(G'\), \(\Delta_+'\), and \(\chi\) agree.

\section{Finite verification data}\label{app:supplement}

The interval input is supplied separately as Supplement S1, version 1.0.0.
The archive name is
\begin{center}\texttt{S1-six-cycle-v1.0.0.zip}\end{center}
Its top-level directory
contains the following files:
\begin{itemize}
\item \texttt{six-cycle-001.json};
\item \texttt{certify\_six\_cycle\_arb.py};
\item \texttt{certify\_candidate\_arb.py};
\item the directed report \texttt{six-cycle-480bit.json}.
\end{itemize}
The reference command is
\begin{verbatim}
python3 certify_six_cycle_arb.py six-cycle-001.json \
  --precision 480 --radius 1e-25 --slices 512 \
  --output certificate.json
\end{verbatim}
Run this command inside the extracted top-level directory.  The reference
dependencies are \texttt{python-flint 0.9.0} and \texttt{mpmath 1.3.0};
S1's \texttt{ENVIRONMENT.md} and \texttt{requirements.txt} give the
environment details.  The command writes \texttt{certificate.json} and
leaves the supplied reference report unchanged.
All decimal inputs are interpreted as exact rational numbers.  Approximate
calculations supply the centres and fixed preconditioners, while directed
interval arithmetic verifies the inclusion, contraction, first-event and
multiplier inequalities.

\section{The six orbits at three scales}\label{app:portraits}

Figure~\ref{fig:six-atlas} shows the six cycles at three spatial scales,
reconstructed from the affine flight data.  Each panel uses equal scales
on the coordinate axes, and all cycles retain their original physical
coordinates.

\begin{figure}[htbp]
\centering
\includerewritefigure{six-cycle-phase-atlas.pdf}
\caption{Numerical reconstruction of the six certified cycles from the
affine flight data.  The panels display successively larger physical
scales; attracting and repelling cycles use solid blue and dashed orange
curves.}
\label{fig:six-atlas}
\end{figure}

\section*{Declaration on AI-assisted tools}

The author used OpenAI Codex for research and writing assistance.
The author reviewed all content,
made the final editorial decisions, and takes full responsibility for the
content of the article.

\begin{thebibliography}{11}
\raggedright
\bibitem{tonnelier}
A. Tonnelier, The McKean's caricature of the FitzHugh--Nagumo model I:
The space-clamped system, \emph{SIAM J. Appl. Math.} \textbf{63} (2002),
459--484.

\bibitem{twofold}
H. Chen, Z. Feng, M. Jia, and Y. Meng, Maximum number of limit cycles of
the piecewise linear Li\'enard system with three zones,
\emph{J. Lond. Math. Soc.} \textbf{113} (2026), e70430.

\bibitem{clzz}
H. Chen, Z. Li, R. Zhang, and X. Zhang, The number of limit cycles of
piecewise linear Li\'enard systems, preprint (2026),
\href{https://arxiv.org/abs/2608.19542v1}{arXiv:2608.19542v1}.

\bibitem{ctz}
H. Chen, Y. Tang, and W. Zhang, Limit cycles in a rotated family of
generalized Li\'enard systems allowing for finitely many switching lines,
\emph{Proc. Roy. Soc. Edinburgh Sect. A} \textbf{155} (2025), 2121--2149.

\bibitem{panazzolo}
D. Panazzolo, Tame semiflows for piecewise linear vector fields,
\emph{Ann. Inst. Fourier} \textbf{52} (2002), 1593--1628.

\bibitem{carmona}
V. Carmona, F. Fern\'andez-S\'anchez, and D. D. Novaes, Uniform upper
bound for the number of limit cycles of planar piecewise linear
differential systems with two zones separated by a straight line,
\emph{Applied Mathematics Letters} \textbf{137} (2023), 108501;
\href{https://arxiv.org/abs/2210.12125v2}{arXiv:2210.12125v2}.

\bibitem{bravo}
J. L. Bravo and R. Trinidad-Forte, Finitude of limit cycles of linear
piecewise ODEs in the cylinder, arXiv:2605.05805v1 [math.CA] (2026).

\bibitem{dmm}
L. van den Dries, A. Macintyre, and D. Marker, The elementary theory of
restricted analytic fields with exponentiation, \emph{Ann. of Math.}
\textbf{140} (1994), 183--205.

\bibitem{vdm}
L. van den Dries and C. Miller, Geometric categories and o-minimal
structures, \emph{Duke Math. J.} \textbf{84} (1996), 497--540.
For Sections 4.2 and 4.4 we use the
\href{https://people.math.osu.edu/miller.1987/newg.pdf}{authors' revised text}
of February 20, 2001.

\bibitem{rump}
S. M. Rump, Verification methods: Rigorous results using floating-point
arithmetic, \emph{Acta Numer.} \textbf{19} (2010), 287--449.

\bibitem{arb}
F. Johansson, Arb: Efficient arbitrary-precision midpoint-radius interval
arithmetic, \emph{IEEE Trans. Comput.} \textbf{66} (2017), 1281--1292.
\end{thebibliography}
\end{document}
